\section{Joint activation and weight zeroing}
\label{sec:joint-zeroing}
We evaluated three combinations on Qwen3-30B-A3B-Instruct-2507 using the same
benchmark-specific generation and grading protocols as the unmodified baseline.
The activation-only arm targets the smallest 50\% of MLP intermediate values
and 25\% of block residual-stream values by magnitude. The weight-only arm
targets the smallest 70\% of input embedding weights and 25\% of both expert
MLP and attention projection weights by magnitude. The third arm
applies both combinations. Their registry names are
\texttt{joint\_activations}, \texttt{joint\_weights}, and \texttt{joint\_all};
\texttt{run.py:JOINT\_ARMS} specifies the intervention strengths and selected
matrices. Weight selection is independent within each output row; activation
selection is per executed token at each selected site. We apply these changes
without fine-tuning. The expert router, output head, and normalization parameters
are preserved.

All 30 full evaluations completed. Before interpreting the scores, we verified
saved configurations, identical baseline question sets and prompts, dataset
manifests, and metrics recomputed from saved samples. Per-site counters confirm
both activation interventions in all 48 layers. Per-matrix reports confirm the
requested weight transformation on the input embedding, 18,432 expert projection
matrices, and 192 attention projection matrices. MMLU-Pro is outside this series.
The question counts are 304 for MMLU STEM, 1,319 for GSM8K, 500 for MATH500,
541 for IFEval, 198 for GPQA Diamond, 164 for HumanEval+, 378 for MBPP+, and
1,300 for each RULER context length.

\begin{table}[ht]
\centering
\caption{Joint-zeroing benchmark scores (\%). IFEval uses strict prompt accuracy;
HumanEval+ and MBPP+ use pass@1 on the extended tests. RULER reports mean
sample scores, including partial credit. Each column uses its own baseline.}
\label{tab:joint-zeroing}
% Hand-authored table transcribed from validated joint-zeroing/full-summary.json.
\resizebox{\textwidth}{!}{%
\begin{tabular}{lrrrrrrrrrr}
\toprule
Experiment & MMLU STEM & GSM8K & MATH500 & IFEval & GPQA & HumanEval+ & MBPP+ & RULER 8k & RULER 16k & RULER 32k \\
\midrule
Baseline & 92.76 & 95.15 & 89.20 & 83.55 & 44.95 & 86.59 & 75.66 & 96.04 & 95.13 & 94.13 \\
Joint activations & 90.13 & 95.07 & 89.60 & 82.44 & 47.98 & 78.66 & 78.57 & 95.91 & 94.51 & 94.05 \\
Joint weights & 91.78 & 95.07 & 84.20 & 84.10 & 37.37 & 87.80 & 78.31 & 95.61 & 94.79 & 94.33 \\
Joint activations + weights & 90.79 & 94.31 & 86.20 & 83.55 & 40.91 & 79.88 & 78.04 & 95.62 & 94.53 & 94.01 \\
\bottomrule
\end{tabular}}

\end{table}

Table~\ref{tab:joint-zeroing} shows task-dependent interactions.
Combining activations preserved MATH500 (89.60\% versus 89.20\% baseline) and
scored 47.98\% on GPQA (44.95\% baseline), but reduced HumanEval+ to 78.66\%
from 86.59\%. The individual MLP and block activation interventions had scored
89.63\% and 85.37\% on HumanEval+, respectively: the joint loss was absent
from either individual intervention in those runs.

Combining weights reduced MATH500 to 84.20\% and GPQA to 37.37\%, below each
individual weight intervention. The individual embedding, MLP, and attention
arms scored 87.00\%, 87.60\%, and 86.60\% on MATH500 and 44.44\%, 44.44\%,
and 42.42\% on GPQA. Adding joint activation zeroing to the joint weight arm
partly recovered these scores to 86.20\% and 40.91\%, while reducing HumanEval+
from 87.80\% to 79.88\%. Thus the measured effects do not add uniformly across
benchmarks. All three combinations stayed within 0.63 percentage points of their
RULER baselines, and the full combination matched baseline IFEval at 83.55\%.

These are single-run comparisons. The baseline repeat increased HumanEval+ by
2.44 percentage points and GPQA by 1.52 points, so apparent improvements require
replicated evaluation before being interpreted as reliable gains. The serving
implementation uses dense kernels; these quality results do not demonstrate a
speedup or memory reduction.

\paragraph{Result provenance.}
{\raggedright
The jobs used the committed source snapshot
\texttt{7831187ddf61f01f9af01d76e26e786528b64372}.
The external experiment artifact directory's \texttt{joint-zeroing/} contains
\path{full-manifest.json} with complete job identifiers,
\path{full-validation.json}, \path{full-summary.json},
\path{full-summary.csv}, and \path{individual-comparisons.json}.
These record all completed evaluations and the individual-arm comparisons used
above. The table transcribes the validated summary, rounded to two decimal places.
\par}
